\documentclass[10pt,a4paper]{amsart}
\usepackage[margin=19mm]{geometry}
\usepackage{amsmath,amssymb,mathtools}
\usepackage[scr=rsfs,cal=euler]{mathalfa}
\usepackage{xfrac}
\usepackage[unicode,hidelinks,bookmarks=false]{hyperref}
\newcommand{\ie}{i.e.\ }
\newcommand{\cf}{cf.\ }
\newcommand{\resp}{respectively\ }
\newcommand{\Subsection}{\subsection}
\providecommand{\nobreakdash}{\nobreak\hbox{-}}
\setlength{\emergencystretch}{2em}
\multlinegap0pt
\usepackage{bm}
\usepackage{bbm}
\usepackage{dsfont}
\usepackage{xspace}
\usepackage{cancel}
\usepackage{mathtools}
\usepackage{cleveref}
\usepackage{amsthm}
\makeatletter
\@ifundefined{theorem}{\newtheorem{theorem}{Theorem}}{}
\@ifundefined{clist}{\newtheorem{clist}{Clist}}{}
\@ifundefined{lemma}{\newtheorem{lemma}[theorem]{Lemma}}{}
\@ifundefined{thmx}{\newtheorem{thmx}{Theorem}}{}
\makeatother
\DeclareMathOperator{\Aut}{Aut}
\DeclareMathOperator{\HT}{HT}
\DeclareMathOperator{\IM}{im}
\newcommand{\Z}{\mathbb{Z}}
\newcommand{\cy}[1]{\Z/{#1}}
\renewcommand{\mod}{\,\,\text{\normalfont mod}\,}
\title[Four-manifolds, two-complexes and the quadratic bias invaria]{Four-manifolds, two-complexes and the quadratic bias invariant}
\author{Ian Hambleton, John Nicholson}
\date{Source: arXiv:2412.15089v3 (CC BY 4.0)}
\begin{document}
\maketitle

\noindent\emph{Source and licence.} Four-manifolds, two-complexes and the quadratic bias invariant. Version arXiv:2412.15089v3 (CC BY 4.0), distributed under the Creative Commons Attribution 4.0 International licence, as recorded in the arXiv licence metadata for this exact version and shown on the abstract page https://arxiv.org/abs/2412.15089v3. Full rights notice: licenses/arxiv-2412.15089-CC-BY-4.0.txt.\par\smallskip

\noindent\textit{Editorial scope.} Counted unit: the Lemma whose source label is \texttt{lemma:Q8-presentations} in the article (source0.tex lines 2049--2059, source label \texttt{lemma:Q8-presentations}), delivered together with the article's own complete original proof of it (source0.tex lines 2061--2081, one \texttt{proof} environment of 21 lines). No mathematics is written, repaired or re-derived: statement and proof are the article's own text line for line. Required context retained uncounted: thmx:main-B\_Q(G) (lines 285--290); lemma:H\_n(G) (lines 1958--1967). Every definition, display and label of those lines is retained unchanged. Omitted from this package and not redistributed: 1--284, 291--1957, 1968--2048, 2060--2060, 2082--2952. Nothing inside a retained excerpt is omitted or altered: every retained body is delivered exactly as the article's lines are written, so no omission receipt of any kind accompanies this package. Citations: the retained excerpts carry no citation command at all, so no bibliography entry of the article is transcribed and no reference is redistributed. Reference adaptation: none was needed and none was made; every reference inside the delivered unit resolves inside it, so no unresolved reference or citation is produced. Printed slips kept literal: no printed word, symbol, hypothesis or number of the counted statement or of its proof was altered. Layout and class: the article's own class is \texttt{amsart} with packages including geometry, xy, appendix, amssymb, amsmath, amsthm, tikz, tikz-cd, mathtools, framed, comment, hyperref, cleveref, rotating; the wrapper is the lane's pinned \texttt{amsart} editorial wrapper at 10pt on A4, which restates the theorem-like environments the retained text needs and adds the article's own macro definitions that the retained text needs (\texttt{\textbackslash Aut}, \texttt{\textbackslash HT}, \texttt{\textbackslash IM}, \texttt{\textbackslash Z}, \texttt{\textbackslash cy}, \texttt{\textbackslash mod}), with extra wrapper packages (bm, bbm, dsfont, xspace, cancel, mathtools, cleveref). The delivered numbering is the unit's own. Assets: no retained span contains a figure environment, a graphics-inclusion command, a PDF-page inclusion, a table or a TikZ picture; no image file is copied into this package, the package declares no asset dependency and no image file is redistributed with it. Licence: version arXiv:2412.15089v3 of this article is distributed under the Creative Commons Attribution 4.0 International licence, as recorded in the arXiv licence metadata for this exact version and shown on the abstract page https://arxiv.org/abs/2412.15089v3. A full rights notice is delivered at licenses/arxiv-2412.15089-CC-BY-4.0.txt. Characters: every retained excerpt is pure ASCII, so the delivered engine, XeLaTeX with its default Latin Modern text font, reports no missing character.
\begin{thmx} \label{thmx:main-B_Q(G)}
Let $G$ be a finite group such that $H_2(G;\Z) \cong (\Z/m)^d$ 
for some $m\geq 1$, $d \ge 3$.  If $G$ is efficient, then there is an isomorphism
\[ B_Q(G) \cong \frac{(\Z/m)^\times}{\pm (\cy m)^{\times 2} \cdot D(G)}\]
where $D(G) = \IM(\varphi_G\colon  \Aut(G) \to  (\Z/m)^\times/\{\pm 1\})$. If $G$ is not efficient, then $B_Q(G)=0$.
\end{thmx}

\begin{lemma} \label{lemma:H_n(G)}
Let $n \ge 2$ and let $G = (\Z / m)^d \times \Z /t $ where $d \geq 3$ and $m, t \ge 1$ are such that $t \mid m$ and $m$, $m/t$ have the same prime factors.
Then
\[
H_n(G) \cong \begin{cases} (\Z/m)^r, & \text{if $n$ is even} \\
 (\Z/m)^r \times \Z/t, & \text{if $n$ is odd}	
 \end{cases}
\]
where $r = r_{(G,n)} \geq 3$.
\end{lemma}

\begin{lemma} \label{lemma:Q8-presentations}
Let $G = Q_8 \times (\Z/p)^3$ where $p$ is a prime such that $p \equiv 1 \mod 4$. Then:
\begin{clist}{(i)}
\item
For each $r \in \Z$ with $(r,p)=1$, we have that
\[ \mathcal{P}_r = \langle A, B, C \mid A^{2p}B^{-2p}, BAB^{-1}A^{2p-1}, C^p, [A,B^{p-1}], [A,C^r], [B,C]\rangle\]
is a presentation for $G$. Furthermore, we can identify $A=xa$, $B=yb$ and $C=c$.
\item
$G$ satisfies the conditions of \cref{thmx:main-B_Q(G)}. More specifically, $G$ satisfies the strong minimality hypothesis, $H_2(G) \cong (\Z/p)^3$, $m_G = p$ and $X_r := X_{\mathcal{P}_r} \in \HT_{\min}(G)$ for each $r$.
\end{clist}
\end{lemma}

\begin{proof}
(i)
Let $(\Z/p)^2 =\langle a, b \rangle$ and $H = Q_8 \times (\Z/p)^2$. Let $A=xa$, $B=yb$. Then $A^p=x$, $A^{1-p}=a$, $B^p=y$ and $B^{1-p}=b$, and so $H = \langle A, B \rangle$. By combining the given presentations for $Q_8$ and $(\Z/p)^2$ together, and adding commutators, we obtain a presentation
\[ \langle A, B \mid A^{2p}B^{-2p}, B^pA^pB^{-p}A^p, A^{p(p-1)}, B^{p(p-1)}, [A^{p-1},B^{p-1}],[A^{p-1},B^p], [A^p,B^{p-1}]\rangle  \]
where we can omit the relators $[A^p,A^{p-1}]$ and $[B^p,B^{b-1}]$ since they are trivial in $F(A,B)$.

The relators $A^{2p}B^{-2p}$ and $B^pA^pB^{-p}A^p$ imply $A^{4p}$ and $B^{4p}$. Since $4 \mid p-1$, this implies $A^{p(p-1)}$ and $B^{p(p-1)}$, so these two relators can be omitted.

Using $[A^{p-1},B^{p-1}]$, we can replace $[A^{p-1},B^p]$ with $[A^{p-1},B]$ and $[A^p,B^{p-1}]$ with $[A,B^{p-1}]$. Either of these relators then implies $[A^{p-1},B^{p-1}]$ and so it can be omitted. We can also use these relators to replace $B^pA^pB^{-p}A^p$ with $BAB^{-1}A^{2p-1}$.
We now have a presentation:
\[ \langle A, B \mid A^{2p}B^{-2p}, BAB^{-1}A^{2p-1}, [A^{p-1},B], [A,B^{p-1}]\rangle.  \]

We now claim that $[A^{p-1},B]$ is a consequence of the other three relators, and so can be omitted. First  note that $[A,B^{p-1}]=1$ and $BAB^{-1}=A^{1-2p}$ implies that $B^{p-1} = AB^{p-1}A^{-1} = (AB^{-1}A^{-1})^{1-p}=(B^{-1}A^{2p})^{1-p}$. Since $A^{2p}=B^{2p}$ is central, we can rewrite this as $B^{p-1} = B^{p-1}A^{2p(1-p)}$ and so $A^{2p(p-1)}=1$. Using $BAB^{-1}=A^{1-2p}$ and $A^{2p(p-1)}=1$, we now obtain $BA^{p-1}B^{-1} = A^{(1-2p)(p-1)} = A^{p-1} A^{2p(1-p)} = A^{p-1}$, as required.
Hence we have that 
\[ \langle A, B \mid A^{2p}B^{-2p}, BAB^{-1}A^{2p-1}, [A,B^{p-1}]\rangle \]
is a presentation for $H$. By adding a generator $C$ and relators $C^p$, $[A,C]$ and $[B,C]$, we clearly obtain a presentation for $G = H \times \Z/p$. It is straightforward to see that, if $r \in \Z$ with $(r,p)=1$, then replacing $[A,C]$ with $[A,C^r]=1$ does not change the group. Thus $\mathcal{P}_r$ presents $G$.

(ii) By the K\"{u}nneth formula, we have that
\[ H_2(G) \cong H_2(Q_8) \oplus H_2((\Z/p)^3) \oplus (H_1(Q_8) \otimes_{\Z} H_1((\Z/p)^3)) \cong H_2((\Z/p)^3) \cong (\Z/p)^3 \]
by \cref{lemma:H_n(G)} and the fact that $H_2(Q_8)=0$, and $H_1(Q_8) \cong (\Z/2)^2$ and $H_1((\Z/p)^3) \cong (\Z/p)^3$ have coprime orders. Since $\chi(X_r) = 1-3+6=4=1+d(H_2(G))$, this implies both that $G$ has the strong minimality hypothesis and that $X_r$ is a minimal finite $2$-complex. Since $H_2(G) \cong (\Z/p)^3$, it now follows that $m_G = p$.
\end{proof}

\end{document}
